% Synthetic M3b fixture: named colours from \palette, used both as
% \textcolor{name}{...} and as bare declarations. Invented content and design.
\documentclass{article}

\fonts{ dir = "fonts", body = "Source Serif Pro", sans = "Open Sans" }

\palette{
  primary = #7C3AED,
  ink     = #18181B,
  muted   = #71717A,
  warn    = #B45309,
}

% A pale tint is deliberate here, so the block says so: decorative entries
% are exempt from the contrast pairing warning.
\palette[decorative]{ soft = primary!30 }

\assert{ pages == 1 }

\begin{document}

\section{\textcolor{primary}{A coloured heading}}

Ordinary text stays black. \textcolor{primary}{This clause is primary} and
\textcolor{muted}{this one is muted}, then black resumes.

{\muted A declaration colours the rest of its group, including
\textbf{bold inside the group}.}

\begin{itemize}
  \item \textcolor{warn}{Warning-coloured item} with a black tail.
  \item Short hex works too: \textcolor{ink}{ink}.
  \item A declared decorative mix: \textcolor{soft}{soft}, and an inline
    pale one that earns the pairing warning: \textcolor{warn!40}{a tint}.
\end{itemize}

\end{document}
